\input{preamble}
\input{format}
\input{commands}
\usepackage{amsmath}
\usepackage{xcolor}
\usepackage{float}
\usepackage{algorithm}
\usepackage{algorithmic}
\definecolor{ECNURed}{RGB}{164,31,53}
\usepackage[fontset=ubuntu]{ctex}




\begin{document}
\title{Quantile-based Risk Aware Reinforcement Learning}
\author{Wu Hao}
\maketitle

\section{Background}
\subsection{Markov Decisioin Process}

We consider a stationary Markov decision process $(\mathcal{S},\mathcal{A},P_{\mathcal{S}},P_{\mathcal{R}},\eta)$ specified by finite state and action spaces $\mathcal{S}$ and $\mathcal{A}$, transition kernel $P_{\mathcal{S}}:\mathcal{S}\times\mathcal{A}\rightarrow\mathcal{P}(\mathcal{S})$, reward distribution function $P_{\mathcal{R}}:\mathcal{S}\times\mathcal{A}\rightarrow\mathcal{P}_1(\mathbb{R})$ and a discount factor $\eta\in[0,1)$, where $\mathcal{P}(\mathcal{S})$ is the set of probability distributions over $\mathcal{S}$ and $\mathcal{P}_1(\mathcal{\mathbb{R}})$ is the set of probability distributions over $\mathbb{R}$ (with its usual Borel $\sigma$-algebra) with finite mean.

Given a policy $\pi: \mathcal{S}\rightarrow\mathcal{P}(\mathcal{A})$ and an initial state $s_0 \in \mathcal{S}$, an agent interacting with the environment according to policy $\pi$ generates a sequence of states, actions and rewards $(S_t,A_t,R_t)_{t=0}^{\infty}$, called a \textit{trajectory}, the joint distribution of which is determined by the transition dynamics and reward distributions of the environment, and the policy of the agent. Here we suppose the Markovian policy is parameterized by a vector $\theta$. More precisely, we have:
\begin{itemize}
    \item $S_0=s_0$ and for each $t\geq0$:
    \item $A_t\mid S_{0:t},A_{0:t-1},R_{0:t-1} \sim \pi(\cdot|s_t;\theta)$;
    \item $R_t\mid S_{0:t},A_{0:t},R_{0:t-1} \sim P_{\mathcal{R}}(\cdot|S_t, A_t)$;
    \item $S_{t+1}\mid S_{0:t},A_{0:t},R_{0:t} \sim P_{\mathcal{S}}(\cdot|S_t,A_t)$;
\end{itemize}
The distribution of the trajectory is thus parameterized by the initial state $S_0$ and policy $\pi$. To illustrate this dependency, we use notation $\mathbb{P}_{s_0}^{\pi}$ and $\mathbb{E}_{s_0}^{\pi}$ to denote the probability and expectation operator corresponding to this distribution, and write $P^{\pi}(\cdot|s)$ for the joint distribution of a reward and next-state pair when the current state is $s$.

\subsection{Cumulative Return}
The quality of the agent's performance on the trajectory is quantified by the \textit{discounted return}, or simply the \textit{return}, which is defined as
\begin{align} \label{eq:return}
    Z:=\sum_{t=0}^{\infty}\eta^t R_t,
\end{align}
where $\eta$ is the discounted factor, $R_t$ is the reward given by the environment at timestep $t$. We write the summation from $0$ to $\infty$, which is a more general notation. If the agent's decision horizon is infinite, which corresponds to the \textit{continuing task}, then (\ref{eq:return}) represents this case, while if the decision horizon is finite up to a constant $T$, which corresponds to the \textit{episodic task}, then we stipulate that $R_t=0$ for all $t>T$. Here and throughout the paper, we consider episodic task where the decision horizon is up to a constant $T$.

The return is a random variable, the randomness of which is determined by the random selection of actions according to policy $\pi(\cdot|X_t)$, the randomness in state transitions $P_{\mathcal{S}}(\cdot|S_t,A_t)$, and the randomness in the observed reward $P_{\mathcal{R}}(\cdot|S_t, A_t)$. Then the trajectory generated by following a policy $\pi$ can be defined as $\tau=\{s_0,a_0,r_0,s_1,a_1,r_1,\dots,s_T\}\sim\Pi(\cdot)$, where $x_0$ is the initial state and $\Pi(\tau)=p(s_0)\Pi_{t=0}^{T-1}\pi(a_t|s_t)p^{\pi}(s_{t+1},r_t|s_t,a_t)=p(s_0)\Pi_{t=0}^{T-1}\pi(a_t|s_t)p_{\mathcal{S}}(s_{t+1}|s_t,a_t)p_{\mathcal{R}}(r_t|s_t,a_t)$. The cumulative reward can then be represented as $Z=\sum_{t=0}^{T-1}\eta^tR_{t}=U(\tau)$, which can be viewed as a function of a random trajectory and follows a distribution $F_{Z}(\cdot;\theta)$.

In classical mean-based RL, the objective is to find the optimal $\theta^*$ that maximizes the expected cumulative return, i.e.
\begin{align} \label{eq:mean}
    \max_{\theta \in \Theta}\mathbb{E}_{Z\sim F_{Z}(\cdot;\theta)}[Z]=\mathbb{E}_{\tau\sim \Pi(\cdot;\theta)}[U(\tau)],
\end{align}
where $\Theta$ is the parameter space.

\subsection{Mean-Based RL}
An important class of techniques for solving (1) is the policy gradient algorithms, where the likelihood-ratio method is frequently used to derive unbiased gradient estimators that do not rely on knowledge of the transition probabilities. Specifically, the gradient of the objective function in \ref{eq:return} can be reformulated as
\begin{align}
\nabla_\theta \mathbb{E}[U(\tau)] 
& = \nabla_\theta \int_{\Omega_\tau} U(\tau)\Pi(\tau;\theta)d\tau = \int_{\Omega_\tau} U(\tau) \frac{\nabla_\theta \Pi(\tau;\theta)}{\Pi(\tau;\theta)} \Pi(\tau;\theta)d\tau  \notag \\
&= \mathbb{E}\left[ U(\tau) \sum_{t=0}^{T-1} \nabla_\theta \log \pi(a_t|s_t;\theta) \right], 
\end{align}
where the interchange of the gradient and integral in the second equality can be justified by the dominated convergence theorem. This yields an unbiased estimator for estimating $\nabla_\theta \mathbb{E}[U(\tau)]$:
$$
U(\tau)\sum_{t=0}^{T-1}\nabla_\theta \log \pi(a_t|s_t;\theta).
$$
In addition, by noticing that for $t' < t$, $\mathbb{E}_{\tau\sim\Pi(\cdot;\theta)}[u(s_{t'+1}, a_{t'}, s_{t'}) \nabla_\theta \log \pi(a_t|s_t;\theta)] = 0$, the gradient estimator can be further reduced to 
$$\sum_{t=0}^{T-1}\left(\nabla_\theta \log \pi(a_t|s_t;\theta) \sum_{t'=t}^{T-1}\eta^{t'-t}u(s_{t'+1}, a_{t'}, s_{t'})\right),$$
where $\sum_{t'=t}^{T-1}\eta^{t'-t}u(s_{t'+1}, a_{t'}, s_{t'})$ is the reward-to-go when starting from state $s_t$ and action $a_t$.

\subsection{Quantile-Based RL}
For a given quantile level $\alpha \in (0,1)$, we define the $\alpha$-quantile for the distribution $F_{Z}(\cdot;\theta)$ as 
\begin{align}
    \notag
    \mathcal{Q}(\alpha;\theta) = \inf\{r: F_{Z}(r;\theta)\geq\alpha\}.
\end{align}
We assume that $F_R(r;\theta)$ is continuously differentiable on $\mathbb{R}$, i.e., $F_R(r;\theta) \in C^{1}(\mathbb{R})$, so that the $\alpha$-quantile can be defined as the inverse of the distribution function, namely $\mathcal{Q}(\alpha;\theta) = F_R^{-1}(\alpha;\theta)$. Our objective is to maximize the $\alpha$-quantile of the distribution over a compact convex parameter space $\Theta \subset \mathbb{R}^m$, i.e.,
\begin{align}
    \max_{\theta \in \Theta} \mathcal{Q}(\alpha;\theta) = \max_{\theta \in \Theta} F_R^{-1}(\alpha;\theta).
\end{align}

As in standard policy gradient methods, we consider solving (\ref{eq:quantile_opt}) using a stochastic gradient ascent approach that leverages gradient information. By the definition of $\mathcal{Q}(\alpha;\theta)$, we have
\begin{align} \notag
    F_Z(\mathcal{Q}(\alpha;\theta); \theta) = \alpha.
\end{align}
Hence, an analytical form of its gradient can be obtained by differentiating both sides of the above equation with respect to $\theta$:
\begin{align}
    \nabla_\theta \mathcal{Q}(\alpha;\theta)
    = -\frac{\nabla_\theta F_Z(r;\theta)}{f_Z(r;\theta)}\bigg|_{r = \mathcal{Q}(\alpha;\theta)}.
    \label{eq:quantile_diff}
\end{align}

The most straightforward way to utilize this relationship is to separately construct estimators for the numerator and the denominator of (\ref{eq:quantile_diff}). However, in practice, two major challenges arise:
\begin{enumerate}
    \item The analytic form of density function $f_{Z}(\cdot;\theta)$ is unknown, and usually hard to get.
    \item The gradient used to optimize estimated quantile include the underlying unknown true quantile $\mathcal{Q}(\alpha;\theta)$, which changes as the underlying parameter vector $\theta$ changes.
\end{enumerate}

In the recently proposed QPO framework, \textcolor{red}{Jiang et al.} observed that the density function $f_{G}(\cdot;\theta)$ remains strictly positive. Consequently, they omitted the density term and focused solely on the gradient direction. And the unknown true quantile $\mathcal{Q}(\alpha;\theta)$ was then estimated via gradient descent, specifically as follows.

For a specific quantile level $\alpha$, the quantile $\mathcal{Q}(\alpha;\theta)$ of distribution $F_{Z}(\cdot;\theta)$ is the minimizer of the problem:
\begin{align}
    \notag
    \mathcal{Q}(\alpha;\theta)=\arg\min_{v}\mathbb{E}_{Z\sim F_{Z}(\cdot;\theta)}\left[(\alpha\mathbf{1}\{Z\geq v\}+(1-\alpha)\mathbf{1}\{Z<v\})|Z-v|\right].
\end{align}
This leads us to optimize the quantile loss:
\begin{align}\label{eq:quantile loss}
\mathbb{E}_{Z}\!\left[\left(\alpha\,\mathbf{1}\{\Delta \ge 0\}+(1-\alpha)\,\mathbf{1}\{\Delta<0\}\right)|\Delta|\right],
\qquad \text{where }\;
\Delta &= Z - v .
\end{align}
Suppose expectation and differentiation are exchangeable, given an observed trajectory $\tau_k$, we then get an unbiased estimator of the negative gradient of (\ref{eq:quantile loss}), thus motivates an update of the form
$$
\mathcal{Q}_{k+1}(\alpha;\theta) = \mathcal{Q}_{k}(\alpha;\theta) + \beta_{k}\left(\alpha\mathbf{1}\{\sum_{t=0}^{T-1}\eta^t r_t\geq v\}+(1-\alpha)\mathbf{1}\{\sum_{t=0}^{T-1}\eta^t r_t<v\}\right),
$$
which can be rewritten as:
\begin{align} \label{quantile update}
    \mathcal{Q}_{k+1}=\mathcal{Q}_k + \beta_k(\alpha-\mathbf{1}\{U(\tau_k)\leq \mathcal{Q}_k\},
\end{align}
where $\beta_k$ is the step-size, $\tau_k$ is the trajectory generated by following the current policy $\pi(\cdot|\cdot;\theta_k)$ parameterized by $\theta_k$, and $\mathcal{Q}_k(\alpha;\theta)$ is an estimate of $q(\alpha;\theta_k)$.

On the other hand, we note that a similar likelihood-ratio technique as in \ref{eq:mean} can be applied to derive the gradient $\nabla_\theta F_Z(r;\theta)$, yielding
\begin{align} \notag
\nabla_\theta F_Z(r;\theta) 
&= \nabla_\theta \mathbb{E}[ \mathbf{1}\{ Z \le r \} ] 
= \nabla_\theta \mathbb{E}[ \mathbf{1}\{ U(\tau) \le r \} ] 
= \nabla_\theta \int_{\Omega_\tau} \mathbf{1}\{ U(\tau) \le r \} \Pi(\tau;\theta) d\tau \notag \\
&= \mathbb{E}\!\left[ \mathbf{1}\{ U(\tau) \le r \} \nabla_\theta \log \Pi(\tau;\theta) \right]
= \mathbb{E}\!\left[ \mathbf{1}\{ U(\tau) \le r \} \sum_{t=0}^{T-1} \nabla_\theta \log \pi(a_t|s_t;\theta) \right] \notag \\
&\approx \frac{1}{N} \sum_{n=0}^{N-1} \mathbf{1}\{ U(\tau_n) \le r \} 
\sum_{t=0}^{T-1} \nabla_\theta \log \pi(a_t^n|s_t^n;\theta). \notag
\end{align}
Consequently, an unbiased estimator for $\nabla_{\theta}F_{Z}(r;\theta)$ is given by
\begin{align} \label{eq:quantile_grad}
    \mathbf{1}\{U(\tau)\leq r\}\sum_{t=0}^{T-1}\nabla_{\theta}\log\pi(a_t|s_t;\theta).
\end{align}

%----------------------------------------------------------------------------

\section{Algorithm 1}
In this paper, we consider the setting where the objective is to optimize the expected cumulative return but add a constraint for the quantile. We formulate the optimization problem as:
% \begin{align} 
%     \max_{\theta \in \Theta}E_{\tau\sim\Pi(\cdot;\theta)}[U(\tau)] ~~\text{s.t.}~\mathcal{Q}(\alpha;\theta)>q,
% \end{align}
\textbf{\textcolor{ECNURed}{
\begin{align} \label{object}
    \max_{\theta \in \Theta}E_{\tau\sim\Pi(\cdot;\theta)}[U(\tau)] ~~\text{s.t.}~F_Z(q;\theta)=\mathbb P_\theta(Z\le q)\le \alpha ,
\end{align}
}}

where $q \in \mathbb{R}$ is a pre-specified threshold.

\subsection{Lagrangian Dual Problem}
Solving problem (\ref{object}) is computationally challenging as the constraint is nonconvex in $\theta$. To generate a feasible high-quality solution, we consider the Lagrangian dual problem as:
% \begin{align} 
%     \min_{\lambda\geq0}\mathcal{L}(\lambda):= \min_{\lambda\geq0}\max_{\theta\in\Theta}\left\{\mathbb{E}_{\tau\sim\Pi(\cdot;\theta)}[U(\tau)]+\lambda\{\mathcal{Q}(\alpha;\theta)-q\}\right\}.
% \end{align}
\textbf{\textcolor{ECNURed}{
\begin{align} \label{dual-lag}
    \min_{\lambda\geq0}\mathcal{L}(\lambda):= \min_{\lambda\geq0}\max_{\theta\in\Theta}\left\{\mathbb{E}_{\tau\sim\Pi(\cdot;\theta)}[U(\tau)]-\lambda\{P_\theta(Z\le q) - \alpha\}\right\}.
\end{align}
}}
This dual problem is convex in $\lambda$, which is a scalar, thus classical convex optimization methods can be used to efficiently find the optimal $\lambda^*$. Then, the corresponding $\tilde{\theta}$ defined as
% \begin{align} 
%     \tilde{\theta}=\mathop{\arg\max}\limits_{\theta\in\Theta} \left\{\mathbb{E}_{\tau\sim\Pi(\cdot;\theta)}[U(\tau)]+\lambda^*\{\mathcal{Q}(\alpha;\theta)-q\}\right\} := \mathop{\arg\max}\limits_{\theta\in\Theta}~\mathcal{J}(\theta)
% \end{align}
\textbf{\textcolor{ECNURed}{
\begin{align} \label{eq:lag-dual}
    \tilde{\theta}=\mathop{\arg\max}\limits_{\theta\in\Theta} \left\{\mathbb{E}_{\tau\sim\Pi(\cdot;\theta)}[U(\tau)]-\lambda^*\{P_\theta(Z\le q) - \alpha\}\right\} := \mathop{\arg\max}\limits_{\theta\in\Theta}~\mathcal{J}(\theta)
\end{align}
}}
is the dual optimal solution and satisfies the quantile constraint.

\subsection{Gradient-Based Algorithm}
In this section we propose a two-time scale gradient-based algorithm for solving (\ref{eq:lag-dual}). Given some $\lambda$, take differentiation in $\mathcal{J}(\theta)$ and combine with the results of (\ref{eq:mean}) and (\ref{eq:quantile_grad}), we have:
% \begin{align} 
%     \nabla_\theta\mathcal{J}(\theta)=\mathbb{E}_{\tau\sim\Pi(\cdot;\theta)}\left[\bigg(U(\tau)\textcolor{red}{-}\frac{\lambda\mathbf{1}\{U(\tau)\leq r\}}{f_{Z}(r;\theta)}\bigg|_{r=\mathcal{Q}(\alpha;\theta)} \bigg)\sum_{t=0}^{T-1}\nabla_{\theta}\log\pi(a_t|s_t;\theta)\right].
% \end{align}
\textbf{\textcolor{ECNURed}{
\begin{align} \label{mix_gradient}
    \nabla_\theta\mathcal{J}(\theta)
    &=\mathbb E_{\tau\sim\Pi(\cdot;\theta)}\left[\Bigl(U(\tau)-\lambda\,\mathbf 1\{U(\tau)\le q\}\Bigr)\sum_{t=0}^\infty\nabla_\theta\log\pi(A_t\mid S_t;\theta)\right]\\
    &=\mathbb E_{\tau\sim\Pi(\cdot;\theta)}\left[\sum_{t=0}^\infty\Bigl(\eta^t G_t-\lambda\,\mathbf 1\{U(\tau)\le q\}\Bigr) \nabla_\theta\log\pi(A_t\mid S_t;\theta)\right].
\end{align}
}}
% The computation of (\ref{mix_gradient}) requires evaluating $f_{Z}(\mathcal{Q}(\alpha;\theta);\theta)$. Since $f_Z$ is unknown, we propose to estimate the density function of $Z$ from its quantile $\mathcal{Q}(\alpha;\theta)$, we have:
% \begin{align}\label{quantile-density}
%     f_Z(Q(\alpha;\theta);\theta)=\lim_{\delta\to0}\frac{2\delta}{Q(\alpha+\delta;\theta)-Q(\alpha-\delta;\theta)},
% \end{align}

% where $\mathcal{Q}(\alpha-\delta;\theta)$ and $\mathcal{Q}(\alpha+\delta;\theta)$ are simultaneously estimated with $\mathcal{Q}(\alpha,\theta)$ using (\ref{q_update}), thus yielding the updating formula:
% \begin{align}
%     \theta_{k+1} &\gets \varphi\left(\theta_k+\gamma_kD(\tau_k,\theta_k,\mathcal{Q}_k)\right), \label{theta_update}\\
%     \mathcal{Q}_{k+1}&\gets\mathcal{Q}_k + \beta_k\left(\alpha-\mathbf{1}\{U(\tau_k)\leq \mathcal{Q}_k\}\right), \label{q_update}
% \end{align}
% where 
% \begin{align}
%     D(\tau_k,\theta_k,\mathcal{Q}_k)):=\Big(U(\tau_k)-\frac{\lambda\mathbf{1}\{U(\tau_k)\leq \mathcal{Q}(\alpha;\theta_k)\}}{\hat{f}_Z(\mathcal{Q}(\alpha;\theta_k);\theta_k)} \Big)\sum_{t=0}^{T-1}\nabla_{\theta}\log\pi(a_t|s_t;\theta_k) \notag
% \end{align}
\textbf{\textcolor{ECNURed}{
\begin{align}
    \theta_{k+1} &\gets \varphi\left(\theta_k+\gamma_kD(\tau_k,\theta_k,\mathcal{Q}_k)\right) \notag \\
    \lambda & \gets \varphi\left(\lambda + \epsilon_k(\mathbf 1\{U(\tau_k)\le q\}-\alpha)\right).
\end{align}
where
\begin{align}
    D(\tau_k,\theta_k,\mathcal{Q}_k))
    &:=\Bigl(U(\tau_k)-\lambda_k\,\mathbf 1\{U(\tau_k)\le q\}\Bigr)\sum_{t=0}^\infty\nabla_\theta \log \pi(A_t^k\mid S_t^k;\theta_k) \notag \\
    &= \sum_{t=0}^\infty \Bigl(\eta^t G_t^k-\lambda_k\,\mathbf 1\{U(\tau_k）\le q\}\Bigr) \nabla_\theta \log \pi(A_t^k\mid S_t^k;\theta_k)
\end{align}
}}
is the gradient estimator.

% We alternately update $\lambda$ and the $\theta_k, ~\mathcal{Q}_k$, and also perform gradient descent updates on $\lambda$, but map it to $[0,+\infty)$ to satisfy the dual Lagrangian constraints (\ref{dual-lag}), that is:
% \begin{align}\label{lambda-update}
%     \lambda \gets \varphi\left(\lambda- \epsilon_k(\mathcal{Q}_k-q)\right).
% \end{align}

We summarize our algorithm in Algorithm \ref{alg:qcrl}
\begin{algorithm}[ht]
\caption{Quantile-Constrained Policy Optimization}
\label{alg:qcrl}
\begin{algorithmic}[1]
\STATE \textbf{Input:} Initial policy $\pi(\cdot|\cdot;\theta_1)$, quantile level $\alpha$, threshold $q$, update interval $\mathcal{I}$, total iterations $N$, step-sizes $\{\gamma_k, \beta_k, \epsilon_k\}$.
\STATE \textbf{Initialize:} $\theta \gets \theta_1$, $\mathcal{Q} \gets \mathcal{Q}_1$, $\lambda \gets 0$.
\FOR{$k = 1, 2, \dots, N$}
    \STATE \textbf{Update Dual Variable (Conditional):}
    \IF{$k \pmod{\mathcal{I}} == 0$}
        \STATE Update $\lambda$ based on the constraint violation:
        \begin{equation}
            \lambda \gets \varphi\left(\lambda + \epsilon_k(\mathcal{Q}_k-q)\right) \notag
        \end{equation}
    \ENDIF
    
    \STATE \textbf{Update Primal Variables ($\theta, \mathcal{Q}$):}
    \STATE 1. Sample a trajectory $\tau_k = \{s_0, a_0, r_0, \dots, s_T\}$ following $\pi(\cdot|\cdot;\theta_k)$.
    \STATE 2. Compute cumulative return: $U(\tau_k) = \sum_{t=0}^{T-1} \eta^t r_t$.
    
    \STATE 3. Update quantile estimate:
    \begin{equation}
        \mathcal{Q}_{k+1} \gets \mathcal{Q}_k + \beta_k \left( \alpha - \mathbf{1}\{ U(\tau_k) \leq \mathcal{Q}_k \} \right) \notag
    \end{equation}
    
    \STATE 4. Estimate density $f_Z$ via (\ref{quantile-density}) and compute policy gradient:
    \begin{equation}
        D_k \gets \left( U(\tau_k) - \frac{\lambda}{\hat{f}_Z} \cdot \mathbf{1}\{ U(\tau_k) \leq \mathcal{Q}_k \} \right) \sum_{t=0}^{T-1} \nabla_\theta \log \pi(a_t | s_t; \theta_k) \notag
    \end{equation}
    
    \STATE 5. Update policy parameter:
    \begin{equation}
        \theta_{k+1} \gets \varphi\left(\theta_k + \gamma_k D_k\right) \notag
    \end{equation}
\ENDFOR
\STATE \textbf{Return} final policy parameter $\theta^*$.
\end{algorithmic}
\end{algorithm}
%----------------------------------------------------------------------------
\section{Algorithm 2}
In Algorithm \ref{alg:qcrl}, a policy update can only be performed after the agent has completed a full interaction with the environment. In some long-tailed distributions, generating a complete trajectory is very costly, thus limiting the data efficiency. In this section, we will use the idea of distributional reinforcement learning, employing Quantile Regression Temporal Difference Learning to improve Algorithm\ref{alg:qcrl}.

We start by introducing a critic that can predict the distribution of future returns. Unlike usual critic that only predict the expectation $V(s):=\mathbb{E}[\sum_{t=0}^{\infty}\eta^tR_{t}|S_0=s]$ or $Q(s,a):=\mathbb{E}[\sum_{t=0}^{\infty}\eta^tR_{t}|S_0=s,A_0=a]$. We need a \textit{distributional critic} to predict the whole return distribution. Several methods have investigated critics in distributional scenarios; for example, QR-DQN\textcolor{blue}{(ref)} is a distributional generalization of the value-based DQN method. Based on the critic model provided by QR-DQN, we combine it with Algorithm \ref{alg:qcrl} to provide an Actor-Critic algorithm for distributional reinforcement learning.

\subsection{Qauntile Distribution and Quantile Regression}
In QR-DQN \textcolor{blue}{(ref)}, the author proposed to estimate the return distribution by a \textit{quantile distribution}. Specifically, for $i=1,...,N$, let $\alpha_i=\frac{i}{N}$, let $\psi: \mathcal{X}\times\mathcal{A}\rightarrow\mathbb{R}^{N}$ be some parametric model for estimating quantiles of $Z(s,a)$, and let $\psi_i$ denotes the $i$-th element of $\psi$. We then have a quantile distribution parameterized by $\psi$ to estimate the return distribution:
\begin{align}\label{quantile distribution}
    Z_{\psi}(s,a):=\frac{1}{N}\sum_{i=1}^{N}\delta_{\psi_i(s,a)},
\end{align}
where $\psi_z$ denotes a Dirac distribution at $z\in\mathbb{R}$.

It has been proved that \textcolor{blue}{(Leamma2)}, for equally spaced quantile levels, if we want to approximate the target distribution $Z(s,a)$ using a quantile distribution like (\ref{quantile distribution}) , then the optimal equally spaced quantile levels under the $W_1$ distance are $\hat{\alpha}_i=\frac{\alpha_{i-1}+\alpha_i}{2}$ for $1\leq i\leq N$. We write $\alpha_0=0$ here for simplicity.

By applying quantile regression, we have that the minimizing values of $\{\psi_1,...,\psi_N\}$ for $W_1(Z,Z_{\psi})$ are those who minimize the following objective:
\begin{align}
    \notag
    \sum_{i=1}^N\mathbb{E}_{\hat{Z}\sim Z}[\rho_{\hat{\alpha}_i}(\hat{Z}-\alpha_i)],
\end{align}
where
\begin{align}\label{quantile loss}
    \rho_{\alpha}(u)=u(\alpha-\delta_{\{u<0\}}), ~~\forall u\in \mathbb{R.}
\end{align}

Since quantile loss is not smooth at zero, we consider a modified quantile loss, called the \textit{quantile Huber loss}:
\begin{align} \label{quantile huber loss}
    \rho_{\alpha}^{\kappa}(u)=|\alpha-\delta_{\{u<0\}}|\frac{\mathcal{L}_{\kappa}(u)}{\kappa},
\end{align}
where
\begin{align}
\mathcal{L}_\kappa(u) = 
    \begin{cases}
        \frac{1}{2}u^2, & \text{if } |u| \leq \kappa \\
        \kappa(|u| - \frac{1}{2}\kappa), & \text{otherwise}
    \end{cases}.
\end{align}
We then have a method for finding the best approximation of quantile distribution $Z_{\psi}(s,a)$ to original return distribution $Z(s,a)$ under $W_1$ distance.

Notice in paricular that the minimizer of the proposed quantile Huber loss function (\ref{quantile huber loss}) is no longer an unbiased estimate of the quantiles, but rather an unbiased estimate of expectile. However, to enable optimization using stochastic gradient descent in deep networks, we can still employ this loss function, and the experimental section demonstrates that its estimation performance is still satisfying.

%---------------------------------------
\subsection{Quantile Regression Temporal Difference Learning}
Recall our Monte-Carlo-based method (\ref{quantile update}) for estimating quantile $\mathcal{Q}_k$. In order to update per transition $(s,a,r,s')$. We employ quantile regression temporal difference learning. Recall the standard TD update for evaluating a policy $\pi$:
\begin{align}
    \notag
    Q(s,a) \gets Q(s,a) + \gamma\left(r+Q(s',a')-Q(s,a)\right), \\
    r\sim R(s,a), s'\sim P(\cdot|s,a), a'\sim \pi(\cdot|s'). \notag
\end{align}
TD allows us to update value function with a single unbiased sample following $\pi$. Quantile regression allows us to improve the estimate of some quantile function $\psi(s,a)$ towards a specific target distribution, $Y(s)$, by observing samples $y\sim Y(s)$ and minimizing Equation \ref{quantile loss}.

We have shown that by minimizing quantile loss we can obtain an approximation distribution $Z_{\psi}(s,a)$ with minimal 1-Wasserstein distance from the original $Z(s,a)$. Combine this with distributional Bellman operator we obtain a target distribution for quantile regression, which leads to the quantile regression temporal differnece learning (QRTD) algorithm summarized as:
\begin{align}
    \psi_i(s,a)\gets \psi_i(s,a) + \beta&\left(\hat{\alpha}_i-\mathbf{1}\{r+\eta z'-\psi_i(s,a)<0\}\right), \notag \\
    r\sim R(s,a), ~s'\sim P(\cdot|s&,a), ~a'\sim\pi(\cdot|s'), ~z'\sim Z_{\psi}(s',a'), \notag  
\end{align}
where $Z_{\psi}$ is a quantile distribution as in (\ref{quantile distribution}), and $\psi_i(s,a)$ is the estimated value of $F^{-1}_{Z^{\pi}(s,a)}(\hat{\alpha}_i)$ in state $s$ taking action $a$.

Notice that when we sample $z'$ from $Z_{\psi}(s',a')$, the obtained $z'$ is actually some $\psi_{i}(s',a')$ since Equation \ref{quantile distribution}. In general, if available, it is better to utilize as much unbiased sample as possible to reduce variance and stabilize the learning process. This leads to the final QRTD updating formula by averaging TD loss for all pairs $\left(\psi_i(s,a),\psi_j(s',a')\right)$ in a single update:
\begin{align}
    \psi_i(s,a)\gets \psi_i(s,a) + \beta\frac{1}{N}&\sum_{j=1}^N\left(\hat{\alpha}_j-\mathbf{1}\{r+\eta \psi_j(s',a')-\psi_i(s,a)<0\}\right),  \\
    \notag  r\sim R(s,a), ~&s'\sim P(\cdot|s,a), ~a'\sim\pi(\cdot|s').
\end{align}

%---------------------------------------
\subsection{Distributional Quantile-Constrained Actor-Critic}

Combine Algorithm \ref{alg:qcrl} with Quantile Regression Temporal Difference Learning, we now propose a data-efficient actor-critic algorithm for the probability-constrained problem. In contrast to the Monte-Carlo-based QCPO update, which requires a complete trajectory before one policy update can be performed, the new algorithm makes use of a distributional critic and one-step temporal-difference learning, allowing the critic and the actor to be updated from single transitions sampled either from the environment or from a replay buffer.

Recall that under the probability-based reformulation, our optimization problem is
\begin{align}
    \max_{\theta\in\Theta}~J(\theta):=\mathbb{E}_{\theta}[Z]
    \qquad \text{s.t.} \qquad
    G(\theta):=\mathbb{P}_{\theta}(Z\le q)\le \alpha,
\end{align}
where
\begin{align}
    Z:=\sum_{t=0}^{\infty}\eta^tR_t
\end{align}
is the infinite-horizon discounted return, $q\in\mathbb{R}$ is a prescribed threshold, and $\alpha\in(0,1)$ is the allowed violation probability level. We consider the associated Lagrangian
\begin{align}
    \mathcal{L}(\theta,\lambda)
    :=
    \mathbb{E}_{\theta}[Z]
    -
    \lambda\bigl(\mathbb{P}_{\theta}(Z\le q)-\alpha\bigr),
    \qquad \lambda\ge 0.
\end{align}

In Algorithm \ref{alg:qcrl}, the actor update is based on the whole-trajectory Monte-Carlo estimator
\begin{align}
    D(\tau_k,\theta_k,\lambda_k)
    :=
    \Bigl(
        U(\tau_k)-\lambda_k\mathbf{1}\{U(\tau_k)\le q\}
    \Bigr)
    \sum_{t=0}^{\infty}
    \nabla_{\theta}\log\pi(A_t^k|S_t^k;\theta_k),
\end{align}
where
\begin{align}
    U(\tau_k)=\sum_{t=0}^{\infty}\eta^tR_t^k
\end{align}
is the realized return of trajectory $\tau_k$. Such an update is correct, but it suffers from low data efficiency. A natural idea is then to introduce a critic and replace the whole-trajectory quantity by a local target. \textcolor{ECNURed}{However, the critic output cannot be directly substituted into the above Monte-Carlo estimator, since the critic estimates a \emph{conditional future return object}, whereas $U(\tau_k)$ is the \emph{realized whole return} of the current sampled trajectory. }\textbf{The correct local actor-critic update must therefore be derived by conditional expectation arguments.}

%---------------------------------------------------------
\subsubsection{Local targets via the tower property}

For each time step $t\ge 0$, define the future discounted return-to-go by
\begin{align}
    G_t := \sum_{\ell=t}^{\infty}\eta^{\ell-t}R_\ell.
\end{align}
Then the gradient of the mean objective can be written as
\begin{align}
    \nabla_{\theta}J(\theta)
    =
    \mathbb{E}_{\theta}\Biggl[
        \sum_{t=0}^{\infty}
        \eta^t G_t
        \nabla_{\theta}\log\pi(A_t\mid S_t;\theta)
    \Biggr].
\end{align}
Applying the tower property to each term gives
\begin{align}
    \nabla_{\theta}J(\theta)
    &=
    \sum_{t=0}^{\infty}
    \mathbb{E}_{\theta}\Bigl[
        \eta^t G_t
        \nabla_{\theta}\log\pi(A_t\mid S_t;\theta)
    \Bigr] \label{eq:mean_local_1}\\
    &=
    \sum_{t=0}^{\infty}
    \mathbb{E}_{\theta}\Bigl[
        \mathbb{E}_{\theta}\bigl[
            \eta^t G_t
            \nabla_{\theta}\log\pi(A_t\mid S_t;\theta)
            \,\big|\,
            S_t,A_t
        \bigr]
    \Bigr] \label{eq:mean_local_2}\\
    &=
    \sum_{t=0}^{\infty}
    \mathbb{E}_{\theta}\Bigl[
        \eta^t
        \mathbb{E}_{\theta}\bigl[
            G_t
            \,\big|\,
            S_t,A_t
        \bigr]
        \nabla_{\theta}\log\pi(A_t\mid S_t;\theta)
    \Bigr], \label{eq:mean_local_3}
\end{align}
where in the last step we use that $\nabla_{\theta}\log\pi(A_t\mid S_t;\theta)$ is measurable with respect to $\sigma(S_t,A_t)$. This motivates the local mean target
\begin{align}
    Q_m^{\pi}(s,a)
    :=
    \mathbb{E}^{\pi}\!\left[
        G_t
        \,\middle|\,
        S_t=s,A_t=a
    \right].
\end{align}

Next, consider the gradient of the chance constraint term
\begin{align}
    G(\theta)
    =
    \mathbb{P}_{\theta}(Z\le q).
\end{align}
By the likelihood-ratio identity,
\begin{align}
    \nabla_{\theta}G(\theta)
    =
    \mathbb{E}_{\theta}\Biggl[
        \mathbf{1}\{Z\le q\}
        \sum_{t=0}^{\infty}
        \nabla_{\theta}\log\pi(A_t\mid S_t;\theta)
    \Biggr].
\end{align}
Applying the tower property again yields
\begin{align}
    \nabla_{\theta}G(\theta)
    &=
    \sum_{t=0}^{\infty}
    \mathbb{E}_{\theta}\Bigl[
        \mathbf{1}\{Z\le q\}
        \nabla_{\theta}\log\pi(A_t\mid S_t;\theta)
    \Bigr] \label{eq:risk_local_1}\\
    &=
    \sum_{t=0}^{\infty}
    \mathbb{E}_{\theta}\Bigl[
        \mathbb{E}_{\theta}\bigl[
            \mathbf{1}\{Z\le q\}
            \nabla_{\theta}\log\pi(A_t\mid S_t;\theta)
            \,\big|\,
            \mathcal{H}_t,A_t
        \bigr]
    \Bigr] \label{eq:risk_local_2}\\
    &=
    \sum_{t=0}^{\infty}
    \mathbb{E}_{\theta}\Bigl[
        \nabla_{\theta}\log\pi(A_t\mid S_t;\theta)
        \,
        \mathbb{E}_{\theta}\bigl[
            \mathbf{1}\{Z\le q\}
            \,\big|\,
            \mathcal{H}_t,A_t
        \bigr]
    \Bigr], \label{eq:risk_local_3}
\end{align}
where $\mathcal{H}_t$ denotes the history up to time $t$.

At this point, one might be tempted to localize the chance-constraint term in the same way as the mean term and directly reuse the mean target $Q_m^\pi$. However, this is in general impossible. The quantity $Q_m^\pi(s,a)$ only characterizes the conditional expectation of the future return, whereas the chance constraint depends on a threshold event. \textbf{\textcolor{ECNURed}{
In particular,
\begin{align}
    \mathbb{P}^{\pi}\!\left(
        Z\le q
        \,\middle|\,
        S_t,A_t
    \right)
    =
    \mathbb{E}^{\pi}\!\left[
        \mathbf{1}\{Z\le q\}
        \,\middle|\,
        S_t,A_t
    \right]
    \neq\mathbf{1}\!\left\{
        \mathbb{E}^{\pi}\!\left[
            Z
            \,\middle|\,
            S_t,A_t
        \right]
        \le q
    \right\}.
\end{align}
}}

Therefore, the constraint gradient cannot be obtained by directly substituting the whole-trajectory quantity with the mean critic.

\textbf{\textcolor{ECNURed}{Next, we proceed from three perspectives:
\begin{itemize}
    \item First, quantile violation event $\{Z\leq q\}$ is a trajectory event that depends on whole trajectory $\tau=(S_0,...,S_t,A_t,...)$, thus $(S_t,A_t)$ is not a sufficient local descriptor, and $\{Z \leq q\}$ cannot be convert to a local event like $G_t$  above.
    \item Secondly, after introducing an augmented state $b_t$, the trajectory event $\{Z\leq q\}$ can be convert to a local event $\{G_t\leq b_t\}.$
    \item Finally, the introduce of $b_t$ will not bring extra complexity since $b_t$ is a series that can be online-updated recursively.
\end{itemize}
}}

\textbf{
To derive the correct local risk target, we first convert the trajectory level event $Z\leq q$ into a local event $G_t\leq b_t$.
} Define the prefix discounted return
\begin{align}
    C_t
    :=
    \sum_{k=0}^{t-1}\eta^kR_k,
    \qquad
    C_0=0.
\end{align}
Then the total return can be decomposed as
\begin{align}
    Z = C_t + \eta^t G_t.
\end{align}
Hence, for any threshold $q$,
\begin{align}
    Z\le q
    \quad \Longleftrightarrow \quad
    G_t \le \frac{q-C_t}{\eta^t}.
\end{align}
This motivates the definition of the remaining budget
\begin{align}\label{budget_definition}
    b_t
    :=
    \frac{q-C_t}{\eta^t}.
\end{align}
Under this definition, the global threshold event becomes
\begin{align}
    Z\le q
    \quad \Longleftrightarrow \quad
    G_t\le b_t.
\end{align}

The introduction of $b_t$ is essential. \textbf{Indeed, the event $\{Z\le q\}$ is not determined by the current state-action pair $(S_t,A_t)$ alone, because it depends not only on the future return $G_t$, but also on the already accumulated prefix return $C_t$.} Thus, two trajectories may share the same current pair $(S_t,A_t)$ while having different past returns, in which case the remaining requirement on the future return is different. \textbf{\textcolor{ECNURed}{The variable $b_t$ summarizes exactly the past information relevant to the threshold event and restores a local Markov representation.}}

\textbf{\textcolor{ECNURed}{
Moreover, $\{b_t\}_{t=0}^{\infty}$ is a series that can be updated online and won't introduce much complexity, since
}}
\begin{align}
    C_{t+1}
    =
    C_t+\eta^tR_t,
\end{align}
the budget admits the one-step recursion
\begin{align}\label{budget_update}
    b_{t+1}
    =
    \frac{q-C_{t+1}}{\eta^{t+1}}
    =
    \frac{b_t-R_t}{\eta}.
\end{align}
Therefore,
\begin{align}
    \mathbb{E}_{\theta}\bigl[
        \mathbf{1}\{Z\le q\}
        \,\big|\,
        \mathcal{H}_t,A_t
    \bigr]
    &=
    \mathbb{P}_{\theta}\bigl(
        G_t\le b_t
        \,\big|\,
        \mathcal{H}_t,A_t
    \bigr).
\end{align}
By the Markov property, once the past dependence has been summarized by $b_t$, the remaining uncertainty only depends on the current state-action pair and the threshold parameter. This motivates the local risk target
\begin{align}
    \Psi^{\pi}(s,a,b)
    :=
    \mathbb{P}^{\pi}\!\left(
        G_t\le b
        \,\middle|\,
        S_t=s,A_t=a
    \right).
\end{align}
Evaluated at the current running budget, it follows that
\begin{align}
    \mathbb{E}_{\theta}\bigl[
        \mathbf{1}\{Z\le q\}
        \,\big|\,
        \mathcal{H}_t,A_t
    \bigr]
    =
    \Psi^{\pi_{\theta}}(S_t,A_t,b_t).
\end{align}
Substituting this into \eqref{eq:risk_local_3}, we obtain
\begin{align}
    \nabla_{\theta}G(\theta)
    =
    \sum_{t=0}^{\infty}
    \mathbb{E}_{\theta}\Bigl[
        \nabla_{\theta}\log\pi(A_t\mid S_t;\theta)
        \,
        \Psi^{\pi_{\theta}}(S_t,A_t,b_t)
    \Bigr].
\end{align}

Combining \eqref{eq:mean_local_3} and the above identity, we arrive at the local representation of the Lagrangian gradient:
\begin{align}\label{local_policy_gradient}
    \nabla_{\theta}\mathcal{L}(\theta,\lambda)
    =
    \sum_{t=0}^{\infty}
    \mathbb{E}_{\theta}\Bigl[
        \nabla_{\theta}\log\pi(A_t\mid S_t;\theta)
        \Bigl(
            \eta^t Q_m^{\pi_{\theta}}(S_t,A_t)
            -
            \lambda \Psi^{\pi_{\theta}}(S_t,A_t,b_t)
        \Bigr)
    \Bigr].
\end{align}
Equation \eqref{local_policy_gradient} is the key step that converts the original whole-trajectory QCPO gradient into a sum of local per-time-step contributions. In particular, it shows that DQC-AC is not obtained by directly replacing the realized whole-trajectory return with a critic output. Rather, the reformulation is achieved by applying the tower property to replace the original whole-trajectory objects with their corresponding local conditional targets: the mean term is localized through the conditional expectation $Q_m^\pi$, while the chance-constraint term is localized through the conditional probability target $\Psi^\pi$ evaluated at the running budget $b_t$. Equation \eqref{local_policy_gradient} is therefore a local decomposition of the exact policy gradient. It is this decomposition that later enables a per-transition stochastic actor update, once $Q_m^\pi$ and $\Psi^\pi$ are replaced by critic-based approximations.

\begin{theorem}{Augmentation MDP}
Fix a policy $\pi$ and a threshold $q\in\mathbb{R}$. Let
\begin{align}
    C_t := \sum_{k=0}^{t-1}\eta^k R_k,
    \qquad
    b_t := \frac{q-C_t}{\eta^t},
\end{align}
and define the augmented state
\begin{align}
    X_t := (S_t,b_t).
\end{align}
Then the following statements hold.

\begin{enumerate}
    \item[(i)] In general, the current state-action pair $(S_t,A_t)$ is not a sufficient local descriptor for the threshold event $\{Z\le q\}$. More precisely, in general there does not exist a measurable function $f:\mathcal{S}\times\mathcal{A}\to [0,1]$ such that
    \begin{align}
        \mathbb{P}^{\pi}\!\left(
            Z\le q
            \,\middle|\,
            \mathcal{H}_t,A_t
        \right)
        =
        f(S_t,A_t)
        \qquad \text{a.s.}
    \end{align}

    \item[(ii)] The augmented process $\{X_t\}_{t\ge 0}$ is a controlled Markov process. In particular, for every measurable set $B\subseteq \mathcal{S}\times\mathbb{R}$,
    \begin{align}
        \mathbb{P}^{\pi}\!\left(
            X_{t+1}\in B
            \,\middle|\,
            \mathcal{H}_t,X_t,A_t
        \right)
        =
        \mathbb{P}^{\pi}\!\left(
            X_{t+1}\in B
            \,\middle|\,
            X_t,A_t
        \right)
        \qquad \text{a.s.}
    \end{align}
\end{enumerate}
Consequently, the local risk target can be written as a function of the current augmented state and action only, namely
\begin{align}
    \mathbb{E}^{\pi}\!\left[
        \mathbf{1}\{Z\le q\}
        \,\middle|\,
        \mathcal{H}_t,A_t
    \right]
    =
    \Psi^{\pi}(S_t,A_t,b_t),
\end{align}
where
\begin{align}
    \Psi^{\pi}(s,a,b)
    :=
    \mathbb{P}^{\pi}\!\left(
        G_t\le b
        \,\middle|\,
        S_t=s,A_t=a
    \right).
\end{align}
\end{theorem}
Proof see appendix.
%---------------------------------------------------------
\subsubsection{Distributional critic and local target extraction}

We now explain how to estimate both local targets using a single distributional critic. As in Sections 3.1 and 3.2, we let the critic output a quantile distribution
\begin{align}
    Z_{\psi}(s,a)
    :=
    \frac{1}{N}\sum_{i=1}^{N}\delta_{\psi_i(s,a)},
\end{align}
where $\psi_i(s,a)$ is the $i$-th quantile estimate.

The target object of the critic is the conditional whole future return distribution, here and throughout $\mathcal{L}$ denotes the law of distribution:
\begin{align}
    \mathcal{Z}^{\pi_{\theta}}(s,a)
    :=
    \mathcal{L}\!\left(
        G_t
        \,\middle|\,
        S_t=s,A_t=a,\text{ afterwards follow }\pi_{\theta}
    \right).
\end{align}
This distribution satisfies the distributional Bellman equation
\begin{align}
    \mathcal{T}^{\pi_{\theta}}\mathcal{Z}(s,a)
    \overset{D}{=}
    R(s,a)+\eta \mathcal{Z}(S',A'),
    \qquad
    S'\sim P(\cdot|s,a),\ A'\sim \pi_{\theta}(\cdot|S').
\end{align}
Hence, the critic estimates the \emph{conditional whole future return distribution} from the current state-action pair under the current policy, rather than the realized whole return of the currently sampled trajectory.

Using the learned quantile distribution, the local mean target is estimated by the mean of the critic output:
\begin{align}\label{mean_from_distribution}
    \widehat{Q}_m(s,a)
    :=
    \mathbb{E}_{Z\sim Z_{\psi}(s,a)}[Z]
    \approx
    \frac{1}{N}\sum_{i=1}^{N}\psi_i(s,a).
\end{align}
The local risk target is obtained by querying the CDF of the same distribution at the current budget:
\begin{align}\label{risk_query}
    \widehat{\Psi}(s,a,b)
    :=
    \widehat{F}_{Z_{\psi}(s,a)}(b)
    :=
    \frac{1}{N}\sum_{i=1}^{N}\mathbf{1}\{\psi_i(s,a)\le b\}.
\end{align}
Therefore, a single distributional critic provides both the local mean signal and the local risk signal required by the actor update.

%---------------------------------------------------------
\subsubsection{Critic update}

Given a sampled transition $(s,a,r,s')$, we sample the next action
\begin{align}
    a'\sim \pi(\cdot|s';\theta),
\end{align}
and define the distributional Bellman targets
\begin{align}
    y_j:=r+\eta \bar{\psi}_j(s',a'),
    \qquad j=1,\dots,N,
\end{align}
where $\bar{\psi}$ denotes the target critic network. The critic parameters are then updated by minimizing the quantile Huber loss
\begin{align}
    \mathcal{L}_{\mathrm{critic}}(\psi)
    :=
    \sum_{i=1}^{N}\sum_{j=1}^{N}
    \rho_{\hat{\alpha}_i}^{\kappa}\bigl(
        y_j-\psi_i(s,a)
    \bigr).
\end{align}
This is a one-step temporal-difference update for the conditional whole future return distribution $\mathcal{Z}^{\pi_{\theta}}(s,a)$. If replay buffer is used, the mathematical object being estimated by the critic does not change; replay buffer only provides more efficient samples for approximating the same Bellman fixed point.

%---------------------------------------------------------
\subsubsection{Actor update}

Using (\ref{mean_from_distribution}) and (\ref{risk_query}), we define the local baselines
\begin{align}
    \widehat{V}_m(s)
    &:=
    \mathbb{E}_{a\sim\pi(\cdot|s;\theta)}[\widehat{Q}_m(s,a)],\\
    \widehat{V}_c(s,b)
    &:=
    \mathbb{E}_{a\sim\pi(\cdot|s;\theta)}[\widehat{\Psi}(s,a,b)].
\end{align}
Then the corresponding local advantages are
\begin{align}
    \widehat{A}_m(s,a)
    &:=
    \widehat{Q}_m(s,a)-\widehat{V}_m(s),\\
    \widehat{A}_c(s,a,b)
    &:=
    \widehat{\Psi}(s,a,b)-\widehat{V}_c(s,b).
\end{align}

To make implementation more convenient, in addition to the running budget $b_t$, we also keep track of the running discount weight
\begin{align}
    d_t:=\eta^t,
    \qquad d_0=1,
    \qquad d_{t+1}=\eta d_t.
\end{align}
Then the local stochastic actor gradient can be written as
\begin{align}
    g_t
    :=
    \nabla_{\theta}\log\pi(a_t|s_t;\theta)
    \Bigl(
        d_t\widehat{A}_m(s_t,a_t)
        -
        \lambda \widehat{A}_c(s_t,a_t,b_t)
    \Bigr).
\end{align}
Accordingly, the actor update is
\begin{align}
    \theta_{k+1}
    \gets
    \varphi\Bigl(
        \theta_k+\gamma_k \widehat{g}_k
    \Bigr),
\end{align}
where $\widehat{g}_k$ is the minibatch average of the local stochastic gradients.

%---------------------------------------------------------
\subsubsection{Dual update}

Finally, note that the global violation probability admits the decomposition
\begin{align}
    G(\theta)
    =
    \mathbb{P}_{\theta}(Z\le q)
    =
    \mathbb{E}_{s_0\sim \rho,\ a_0\sim \pi(\cdot|s_0;\theta)}
    \Bigl[
        F_{\mathcal{Z}^{\pi_{\theta}}(s_0,a_0)}(q)
    \Bigr].
\end{align}
Using the critic, we estimate it by
\begin{align}
    \widehat{G}(\theta)
    :=
    \mathbb{E}_{s_0\sim \rho,\ a_0\sim \pi(\cdot|s_0;\theta)}
    \bigl[
        \widehat{\Psi}(s_0,a_0,q)
    \bigr].
\end{align}
Then the dual variable is updated by projected gradient ascent:
\begin{align}
    \lambda_{k+1}
    \gets
    \varphi\Bigl(
        \lambda_k+\epsilon_k\bigl(\widehat{G}(\theta_k)-\alpha\bigr)
    \Bigr),
\end{align}
where here $\varphi(\cdot)$ denotes projection onto $[0,\infty)$.

%---------------------------------------------------------
\subsubsection{Algorithmic Procedure}

We summarize the resulting DQC-AC algorithm below.

\begin{algorithm}[ht]
\small % 缩小字号以适应页面
\caption{Distributional Quantile-Constrained Actor-Critic (DQC-AC)}
\label{alg:dqcac}
\begin{algorithmic}[1]
\STATE \textbf{Input:} $\pi(\cdot|\cdot;\theta_1)$, $\psi(\cdot|\cdot;\omega_1)$ ($N$ quantiles), $q$, $\alpha$, $\mathcal{I}$, $K$, step-sizes $\{\gamma_k,\epsilon_k,\xi_k\}$.
\STATE \textbf{Initialize:} $\theta\gets \theta_1, \omega\gets \omega_1, \bar{\omega}\gets \omega_1, \mathcal{D}\gets \emptyset, \lambda\gets 0$.
\FOR{$k=1,2,\dots,K$}
    \STATE \textit{\textbf{Collect rollout:}} Reset env, $b_0\gets q, d_0\gets 1$.
    \FOR{$t=0, 1, \dots$}
        \STATE Sample $a_t\sim \pi(\cdot|s_t;\theta_k)$, observe $(r_t,s_{t+1})$, store $(s_t,a_t,r_t,s_{t+1},b_t,d_t)$ in $\mathcal{D}$.
        \STATE $b_{t+1} \gets (b_t-r_t)/\eta, \quad d_{t+1} \gets \eta d_t$.
    \ENDFOR

    \STATE \textit{\textbf{Critic Update:}} Sample minibatch $\{(s_m,a_m,r_m,s'_m,b_m,d_m)\}_{m=1}^{M}$.
    \STATE $y_{m,j}=r_m+\eta \bar{\psi}_j(s'_m,a'_m)$ for $j=1,\dots,N$; Update $\omega$ by minimizing:
    \STATE $\omega_{k+1} \gets \omega_k-\xi_k\nabla_{\omega} \frac{1}{M}\sum_{m, i, j} \rho_{\hat{\alpha}_i}^{\kappa}(y_{m,j}-\psi_i(s_m,a_m))$.

    \STATE \textit{\textbf{Actor Update:}} For each $m$, compute:
    \STATE $\widehat{Q}_m = \frac{1}{N}\sum_i\psi_i(s_m,a_m)$, $\widehat{\Psi}_m = \frac{1}{N}\sum_i\mathbf{1}\{\psi_i(s_m,a_m)\le b_m\}$.
    \STATE $\widehat{V}_m(s_m) \approx \widehat{Q}_m(s_m,\tilde{a}_m)$, $\widehat{V}_c(s_m,b_m) \approx \widehat{\Psi}(s_m,\tilde{a}_m,b_m)$ where $\tilde{a}_m \sim \pi(\cdot|s_m;\theta_k)$.
    \STATE $g_m = \nabla_{\theta}\log\pi(a_m|s_m;\theta_k) \bigl[ d_m(\widehat{Q}_m-\widehat{V}_m) - \lambda_k(\widehat{\Psi}_m-\widehat{V}_c) \bigr]$.
    \STATE $\theta_{k+1} \gets \varphi( \theta_k+\gamma_k\frac{1}{M}\sum_{m=1}^{M}g_m )$.

    \STATE \textit{\textbf{Dual Update (Conditional):}} \IF{$k \pmod{\mathcal{I}} == 0$}
        \STATE Sample $\{s_0^{(m)}\}$, compute $\widehat{G}(\theta_{k+1}) = \frac{1}{M_0 N}\sum_{m,i} \mathbf{1}\{\psi_i(s_0^{(m)},a_0^{(m)})\le q\}$.
        \STATE $\lambda_{k+1} \gets \varphi\bigl( \lambda_k+\epsilon_k(\widehat{G}(\theta_{k+1})-\alpha) \bigr)$.
    \ENDIF
    \STATE Periodically update target critic $\bar{\omega}$.
\ENDFOR
\STATE \textbf{Return} $\theta^*$.
\end{algorithmic}
\end{algorithm}

\paragraph{Remark.}
The critic in Algorithm \ref{alg:dqcac} estimates the conditional whole future return distribution $\mathcal{Z}^{\pi_{\theta}}(s,a)$, not the realized whole return of the currently sampled trajectory. The role of the budget variable $b_t$ is to convert the global threshold event $\{Z\le q\}$ into the local event $\{G_t\le b_t\}$. Therefore, the local risk target is obtained by querying the CDF of the critic output at the current budget, rather than by directly replacing the whole-trajectory return in the original QCPO gradient estimator.

\newpage

\section{Appendix}
\subsection{Proof of Theorem 1}
\begin{proof}
\medskip
\noindent\textbf{Proof of (i).}
Recall that
\begin{align}
    Z = C_t + \eta^t G_t,
\end{align}
and therefore
\begin{align}
    Z\le q
    \quad \Longleftrightarrow \quad
    G_t \le \frac{q-C_t}{\eta^t}
    =
    b_t.
\end{align}
Hence,
\begin{align}
    \mathbb{P}^{\pi}\!\left(
        Z\le q
        \,\middle|\,
        \mathcal{H}_t,A_t
    \right)
    =
    \mathbb{P}^{\pi}\!\left(
        G_t\le b_t
        \,\middle|\,
        \mathcal{H}_t,A_t
    \right).
\end{align}
Now suppose two histories $h_t$ and $h_t'$ satisfy
\begin{align}
    S_t(h_t)=S_t(h_t')=s,
    \qquad
    A_t(h_t)=A_t(h_t')=a,
\end{align}
but have different prefix returns
\begin{align}
    C_t(h_t)\neq C_t(h_t').
\end{align}
Equivalently, their induced budgets are different:
\begin{align}
    b_t(h_t)\neq b_t(h_t').
\end{align}
Let
\begin{align}
    F_{s,a}^{\pi}(x)
    :=
    \mathbb{P}^{\pi}\!\left(
        G_t\le x
        \,\middle|\,
        S_t=s,A_t=a
    \right)
\end{align}
denote the conditional distribution function of the future return-to-go. Then
\begin{align}
    \mathbb{P}^{\pi}\!\left(
        Z\le q
        \,\middle|\,
        \mathcal{H}_t=h_t,A_t=a
    \right)
    &=
    F_{s,a}^{\pi}\!\bigl(b_t(h_t)\bigr),\\
    \mathbb{P}^{\pi}\!\left(
        Z\le q
        \,\middle|\,
        \mathcal{H}_t=h_t',A_t=a
    \right)
    &=
    F_{s,a}^{\pi}\!\bigl(b_t(h_t')\bigr).
\end{align}
In general, these two quantities are different whenever $F_{s,a}^{\pi}$ is not constant on the interval between $b_t(h_t)$ and $b_t(h_t')$. Therefore, the conditional violation probability cannot in general be represented as a function of $(S_t,A_t)$ alone. This proves (i).

\medskip
\noindent\textbf{Proof of (ii).}
From the definition of $b_t$, we have the one-step recursion
\begin{align}
    b_{t+1}
    =
    \frac{q-C_{t+1}}{\eta^{t+1}}
    =
    \frac{q-(C_t+\eta^t R_t)}{\eta^{t+1}}
    =
    \frac{b_t-R_t}{\eta}.
\end{align}
Thus, once $(S_t,b_t)$ and $A_t$ are given, the next budget $b_{t+1}$ is completely determined by the one-step reward $R_t$.

Let $\mathsf{P}(dr,ds'\mid s,a)$ denote the joint law of $(R_t,S_{t+1})$ conditional on $(S_t,A_t)=(s,a)$. For any measurable set $B\subseteq \mathcal{S}\times\mathbb{R}$,
\begin{align}
    &\mathbb{P}^{\pi}\!\left(
        X_{t+1}\in B
        \,\middle|\,
        \mathcal{H}_t,X_t=(s,b),A_t=a
    \right) \notag\\
    &\qquad =
    \mathbb{P}^{\pi}\!\left(
        \left(S_{t+1},\frac{b-R_t}{\eta}\right)\in B
        \,\middle|\,
        \mathcal{H}_t,S_t=s,b_t=b,A_t=a
    \right) \notag\\
    &\qquad =
    \int
    \mathbf{1}\!\left\{
        \left(s',\frac{b-r}{\eta}\right)\in B
    \right\}
    \mathsf{P}(dr,ds'\mid s,a),
\end{align}
where the last equality follows from the Markov property of the original MDP: conditioned on the current state-action pair $(S_t,A_t)=(s,a)$, the joint law of $(R_t,S_{t+1})$ is independent of the earlier history. Define
\begin{align}
    K(B\mid (s,b),a)
    :=
    \int
    \mathbf{1}\!\left\{
        \left(s',\frac{b-r}{\eta}\right)\in B
    \right\}
    \mathsf{P}(dr,ds'\mid s,a).
\end{align}
Then
\begin{align}
    \mathbb{P}^{\pi}\!\left(
        X_{t+1}\in B
        \,\middle|\,
        \mathcal{H}_t,X_t,A_t
    \right)
    =
    K(B\mid X_t,A_t)
    \qquad \text{a.s.},
\end{align}
which shows that $\{X_t\}_{t\ge 0}$ is a controlled Markov process.

Finally, since
\begin{align}
    Z\le q
    \quad \Longleftrightarrow \quad
    G_t\le b_t,
\end{align}
we have
\begin{align}
    \mathbb{E}^{\pi}\!\left[
        \mathbf{1}\{Z\le q\}
        \,\middle|\,
        \mathcal{H}_t,A_t
    \right]
    &=
    \mathbb{P}^{\pi}\!\left(
        G_t\le b_t
        \,\middle|\,
        \mathcal{H}_t,A_t
    \right) \\
    &=
    \Psi^{\pi}(S_t,A_t,b_t),
\end{align}
where the second equality follows because, after augmenting the state by $b_t$, the remaining uncertainty is fully characterized by the current augmented state-action triple $(S_t,b_t,A_t)$. This proves the claim.
\end{proof}





%------------------------------------------------------------------------------
%-----------------------------------------------------------------------------
\end{document}

